\begin{supp}[モノイド環]
  群環や多項式環は, モノイド環の特別な場合とみることができる.
  ここではまずモノイド環の定義を述べる. $A$を零環でない可換環, $M$をモノイドとする. モノイド環$A[M]$とは,
  写像$a \colon M \to A$であって, 有限個の$m \in M$を除き$a(m) = 0$となるものの全体
  \[
    A[M] = \left\{ a \colon M \to A \ \middle|\ \text{有限個を除き } a(m) = 0 \right\}
  \]
  に, $A$作用と和と積を次のように定めたものである:
  まず, $A$作用と和を
  \[
    (a + b)(m) := a(m) + b(m), \quad (\lambda a)(m) := \lambda a(m) \quad (\lambda \in A)
  \]
  と定める. ここで$A[M]$の特別な元$e_s$を
  \[
    e_s(m) := \begin{cases}
      1_A & (m = s) \\
      0_A & (m \neq s)
    \end{cases}
  \]
  と定めれば, これを用いて$A[M]$の任意の元$a$は, 上で定めた$A$作用と和を用いて
  \[
    a = \sum_{m \in M} a(m) e_m
  \]
  と一意的に表される. 積は$e_m$たちの積を
  \[
    e_m e_n = e_{mn}
  \]
  として, $A$双線型に拡張することで定める. 具体的には, $a = \sum_{m \in M} a(m) e_m$, $b = \sum_{n \in M} b(n) e_n$に対して
  \[
    ab = \sum_{m,n \in M} a(m) b(n) e_{mn} = \sum_{p \in M}\Biggl(\,\sum_{\substack{m,n \in M \\ mn = p}} a(m) b(n)\Biggr) e_p
  \]
  と定める. 以上のように定めた和と積により, $A[M]$は可換環$A$上の環となる. これをモノイド環という.
  $s \in M$を$e_s \in A[M]$と同一視すれば, $A[M]$の任意の元は
  \[
    a = \sum_{m \in M} a(m) m
  \]
  と表すことができる.

  % 例: 群環
  \begin{ex}
    群$G$をモノイドとみなせば, 群環$A[G]$はモノイド環の特別な場合である.
  \end{ex}
  % 例: 多項式環
  \begin{ex}
    可換環$A$に対して, 多項式環$A[x]$はモノイド$\mathbb{N}$に対するモノイド環$A[\mathbb{N}]$である.
  \end{ex}
\end{supp}

\begin{ans}{1.1.1}
  \begin{align*}
   gh &= (3\sigma_6 - 2\sigma_4 + \sigma_2)
  + (-6\sigma_2\sigma_6 + 4\sigma_2\sigma_4 - 2\sigma_2^2)
  + (9\sigma_5\sigma_6 - 6\sigma_5\sigma_4 + 3\sigma_5\sigma_2) \\
  &= (3\sigma_6 - 2\sigma_4 + \sigma_2)
  + (-6\sigma_3 + 4\sigma_5 - 2)
  + (9 -6\sigma_2 + 3\sigma_3) \\
  &= 7 - 5\sigma_2 - 3\sigma_3 - 2\sigma_4 + 4\sigma_5 + 3\sigma_6
  \end{align*}
\end{ans}

\begin{ans}{1.1.2}
  $G = \{g_1, g_2, \dots, g_n\}$とする.
  $gg_i=g_{j_i}$とおけば, $g_{j_i} (i=1, 2, \dots, n)$は互いに異なる$n$個の元である. したがって,
  \[
    gN = g\sum_{i=1}^n g_i = \sum_{i=1}^n gg_i = \sum_{i=1}^n g_{j_i} = \sum_{i=1}^n g_i = N
  \]
  である.$Ng = N$も同様.
\end{ans}

\begin{ans}{1.1.3}
  略.
\end{ans}
